\section{Open questions and next steps}

Five substantive questions remain open after this replication. Each is
concrete, testable on free endpoints, and directly connected to a gap
either in the paper or in our reproduction envelope.

\subsection*{Q1. Wall-clock vs modern classical Bayesian inference}
\textbf{Basis.} The paper claims a polynomial asymptotic speedup but
never benchmarks against 2018-era or later classical Bayesian
inference. Our replication is a $2\times 2$ primitive check; the
asymptotic claim is untested empirically.
\textbf{Next steps.} Pick a mid-size GP regression benchmark
(UCI protein or elevators, $N \sim 10^4$). Run (a) exact GP with
Cholesky, (b) SVI GP, (c) HMC over hyperparameters + exact posterior,
(d) HHL-style noiseless simulator on the same kernel matrix where
feasible, or a numerical HHL proxy where not. Compare posterior
mean/variance MSE and wall-clock. Free endpoints: CPU / uicgpu A100
for classical; Qiskit-Aer statevector for HHL sim.

\subsection*{Q2. Realistic-hardware noise sweep on current IBM devices}
\textbf{Basis.} C4 ($F=0.78$ on IBMQX5) cannot be reproduced --- device
retired. Our depolarizing sweep at 5--10\% brackets it, but modern IBM
hardware has $\sim 10^{-3}$ per-gate error and different dominant
error channels (crosstalk, readout, coherent drift).
\textbf{Next steps.} Pull the current calibration \texttt{NoiseModel}
from an IBM Quantum backend via \texttt{qiskit-ibm-runtime} (free
tier); rerun the $2\times 2$ HHL circuit; report fidelity vs
classical target. If free-tier allotment permits, submit the actual
4-qubit circuit as a small job to a Heron-family device.

\subsection*{Q3. Uncertainty calibration on real data}
\textbf{Basis.} We reproduce the numerical variance ($0.6725$) but never
test coverage. Neither does the paper. Calibration is the practical
distinguisher between ``Bayesian'' and ``point estimate + error bar'',
and HHL-specific noise (finite-precision QPE, post-selection
dilution) could bias the reported variance.
\textbf{Next steps.} Small regression benchmark (Boston, or synthetic
Snelson 1-D), leave-one-out CV. For each held-out point, compute
HHL-predicted mean and variance under the noisy simulator
(\texttt{gate\_noise} $\in [0.001, 0.05]$); measure empirical
coverage of 68/95/99\% credible intervals; report reliability
diagrams.

\subsection*{Q4. Kernel-choice condition-number scaling and preconditioning}
\textbf{Basis.} HHL runtime scales with $\kappa$, and Aaronson's
fine-print critique of HHL centers on hidden $\kappa$/preparation
costs. The paper does not analyze $\kappa$ on any real dataset.
\textbf{Next steps.} Load 3--5 standard GP-benchmark datasets, build
$K + \sigma_n^2 I$ under RBF/Mat\'ern-3/2/Mat\'ern-5/2 kernels with
marginal-likelihood-fit hyperparameters. Report empirical $\kappa$
distribution. Apply Nystr\"om / SKI / inducing-point preconditioning;
report reduced effective $\kappa$. Fold into the theoretical HHL
runtime formula and compare to classical Cholesky.

\subsection*{Q5. Hybrid quantum-classical active learning for experiment design}
\textbf{Basis.} The paper stops at posterior inference. GP posterior
variance is the standard active-learning acquisition function, and
quantum-experiment design (VQE, quantum optimal control) is precisely
where a quantum-native posterior might most naturally couple.
\textbf{Next steps.} Pick a small VQE benchmark (H$_2$, LiH
ansatz-parameter search) or a QOC pulse-shape task. Build a GP
surrogate over the parameter landscape. Run classical GP active
learning with expected-improvement / max-variance acquisition. Replace
the posterior-variance step with the HHL-based inference from this
replication (noiseless simulator first). Compare queries-to-target
against the classical baseline. Free: local Qiskit-Aer + classical
VQE simulator.
